% !TeX program = xelatex
\documentclass[11pt,a4paper]{ctexart}
% Shared preamble for Agena architecture and process documents.
\usepackage[a4paper,margin=20mm]{geometry}
\usepackage{array}
\usepackage{booktabs}
\usepackage{caption}
\usepackage{enumitem}
\usepackage{fancyhdr}
\usepackage{hyperref}
\usepackage{longtable}
\usepackage{pdflscape}
\usepackage{tabularx}
\usepackage{tikz}
\usepackage[most]{tcolorbox}
\usepackage{xcolor}
\usepackage{listings}
\usepackage{url}

\usetikzlibrary{
  arrows.meta,
  backgrounds,
  calc,
  fit,
  positioning,
  shapes.geometric,
  shapes.misc,
  shapes.multipart
}

\definecolor{agenaBlue}{HTML}{1E5AA8}
\definecolor{agenaBlueSoft}{HTML}{EAF2FD}
\definecolor{agenaGreen}{HTML}{1D7A5E}
\definecolor{agenaGreenSoft}{HTML}{E7F6F1}
\definecolor{agenaOrange}{HTML}{B85C13}
\definecolor{agenaOrangeSoft}{HTML}{FCEFE4}
\definecolor{agenaPurple}{HTML}{5A4FCF}
\definecolor{agenaPurpleSoft}{HTML}{ECEAFB}
\definecolor{agenaInk}{HTML}{1D2433}
\definecolor{agenaMist}{HTML}{F7F9FC}
\definecolor{agenaSand}{HTML}{FFF9F0}
\definecolor{agenaGray}{HTML}{5D6678}
\definecolor{agenaBorder}{HTML}{C8D2E1}

\hypersetup{
  colorlinks=true,
  linkcolor=agenaBlue,
  urlcolor=agenaBlue,
  citecolor=agenaGreen,
  pdftitle={Agena Documentation},
  pdfauthor={OpenAI Codex}
}

\ctexset{
  section = {
    format = \Large\bfseries\color{agenaInk}
  },
  subsection = {
    format = \large\bfseries\color{agenaInk}
  },
  subsubsection = {
    format = \normalsize\bfseries\color{agenaInk}
  }
}

\setlength{\parskip}{0.55em}
\setlength{\parindent}{0pt}
\setlength{\headheight}{14.5pt}
\setlist[itemize]{leftmargin=2.2em}
\setlist[enumerate]{leftmargin=2.4em}
\captionsetup{hypcap=false}

\pagestyle{fancy}
\fancyhf{}
\fancyhead[L]{\textcolor{agenaGray}{Agena 文档}}
\fancyhead[R]{\textcolor{agenaGray}{\leftmark}}
\fancyfoot[C]{\textcolor{agenaGray}{\thepage}}
\renewcommand{\headrulewidth}{0.3pt}
\renewcommand{\footrulewidth}{0pt}

\newcommand{\docversion}{2026-04-06}
\newcommand{\pathinline}[1]{\texttt{\nolinkurl{#1}}}
\newcommand{\cmdinline}[1]{\texttt{#1}}
\newcommand{\term}[1]{\textbf{\textcolor{agenaInk}{#1}}}
\newcommand{\highlight}[1]{\textcolor{agenaBlue}{\textbf{#1}}}

\newtcolorbox{summarybox}[1][]{
  breakable,
  colback=agenaBlueSoft,
  colframe=agenaBlue,
  boxrule=0.8pt,
  arc=2pt,
  left=8pt,
  right=8pt,
  top=6pt,
  bottom=6pt,
  title=#1,
  fonttitle=\bfseries
}

\newtcolorbox{notebox}[1][]{
  breakable,
  colback=agenaSand,
  colframe=agenaOrange,
  boxrule=0.8pt,
  arc=2pt,
  left=8pt,
  right=8pt,
  top=6pt,
  bottom=6pt,
  title=#1,
  fonttitle=\bfseries
}

\newtcolorbox{techbox}[1][]{
  breakable,
  colback=agenaGreenSoft,
  colframe=agenaGreen,
  boxrule=0.8pt,
  arc=2pt,
  left=8pt,
  right=8pt,
  top=6pt,
  bottom=6pt,
  title=#1,
  fonttitle=\bfseries
}

\lstdefinestyle{agenaCode}{
  basicstyle=\ttfamily\small,
  backgroundcolor=\color{agenaMist},
  frame=single,
  rulecolor=\color{agenaBorder},
  columns=fullflexible,
  keepspaces=true,
  breaklines=true,
  showstringspaces=false
}

\tikzset{
  diagram/.style={
    font=\footnotesize,
    line join=round,
    line cap=round
  },
  module/.style={
    rectangle,
    rounded corners=4pt,
    draw=agenaInk!85,
    fill=white,
    thick,
    align=left,
    inner sep=8pt,
    text width=38mm,
    minimum height=12mm
  },
  modulewide/.style={
    module,
    text width=58mm
  },
  terminal/.style={
    module,
    draw=agenaBlue!80!black,
    fill=agenaBlueSoft
  },
  runtime/.style={
    module,
    draw=agenaGreen!80!black,
    fill=agenaGreenSoft
  },
  storage/.style={
    module,
    draw=agenaOrange!85!black,
    fill=agenaOrangeSoft
  },
  external/.style={
    module,
    draw=agenaPurple!80!black,
    fill=agenaPurpleSoft
  },
  note/.style={
    rectangle,
    rounded corners=4pt,
    draw=agenaBorder,
    fill=agenaMist,
    align=left,
    inner sep=6pt,
    text width=34mm,
    font=\scriptsize
  },
  decision/.style={
    diamond,
    draw=agenaInk!85,
    fill=agenaSand,
    aspect=2.2,
    align=center,
    inner sep=2pt,
    text width=24mm
  },
  group/.style={
    draw=agenaBorder,
    rounded corners=6pt,
    fill=agenaMist,
    inner sep=10pt
  },
  flow/.style={
    -{Latex[length=3mm,width=2mm]},
    thick,
    draw=agenaInk!85
  },
  signal/.style={
    flow,
    dashed,
    draw=agenaBlue!80!black
  },
  weakflow/.style={
    -{Latex[length=2.5mm,width=1.8mm]},
    semithick,
    draw=agenaGray
  },
  dbentity/.style={
    rectangle split,
    rectangle split parts=2,
    rectangle split part fill={agenaOrangeSoft,white},
    draw=agenaOrange!85!black,
    rounded corners=3pt,
    thick,
    align=left,
    text width=43mm,
    inner sep=6pt
  }
}

\newenvironment{diagramlandscape}
  {\clearpage\begin{landscape}\thispagestyle{plain}\begin{center}\captionsetup{width=0.94\linewidth}}
  {\end{center}\end{landscape}\clearpage}


\title{\textbf{Agena 流程文档}\\\large 面向小白的端到端执行流程说明}
\author{基于仓库源码整理}
\date{\docversion}

\begin{document}
\maketitle
\tableofcontents
\clearpage

\section{文档目标}

架构文档回答的是“系统由什么组成”；这份流程文档回答的是“\highlight{系统到底怎么动起来}”。

对于第一次接触 Agena 的人，最值得搞懂的不是零碎 API，而是下面这些完整流程：

\begin{enumerate}
  \item 运行时如何启动、如何热重载；
  \item 一次用户提问如何变成一条 assistant 回答；
  \item 模型发起工具调用时，系统如何决定直接执行、申请权限，还是向用户追问；
  \item 前端为什么可以一边看流式输出、一边通过 SSE 增量刷新；
  \item OpenAI / GitLab / Copilot 这些登录流程最后是怎么落到本地 auth store 的。
\end{enumerate}

\begin{summarybox}[读流程文档时的建议]
不要试图一次把所有细节背下来。你只要先抓住“\highlight{主干流程}、\highlight{阻塞点}、\highlight{恢复点}”这三类节点，后面的源码阅读就会顺很多。
\end{summarybox}

\section{先看全局主流程}

先给出一个压缩过的整体图。你可以把它当成“从用户输入到稳定态”的总路线图。

\begin{diagramlandscape}
\begin{tikzpicture}[diagram, node distance=12mm and 16mm]
  \node[terminal, text width=34mm] (user) {\textbf{1. 用户输入}\\来自 ChatPage、HTTP client 或桌面 UI};
  \node[terminal, right=of user, text width=40mm] (endpoint) {\textbf{2. 入口 API}\\\pathinline{POST /sessions/\{id\}/turns}\\或 continue / reply 系列接口};
  \node[runtime, right=of endpoint, text width=42mm] (persistuser) {\textbf{4. 写入用户消息}\\\pathinline{submit_user_turn()}\\先把 user message 落库};
  \node[runtime, right=of persistuser, text width=40mm] (loop) {\textbf{5. run\_until\_stable}\\决定下一步是跑模型、执行工具、还是停下};

  \node[runtime, below=14mm of endpoint, text width=40mm] (versioncheck) {\textbf{3. 版本检查}\\如果请求带 \pathinline{If-Match}\\先校验 session version};
  \node[runtime, below=14mm of persistuser, text width=42mm] (tool) {\textbf{6. 工具分支}\\权限检查、ask\_user、执行工具、生成 tool result message};
  \node[runtime, below=14mm of loop, text width=40mm] (model) {\textbf{7. 模型分支}\\SessionProcessor 打开 provider stream\\MessageProjector 生成 delta};
  \node[runtime, below=14mm of model, text width=42mm] (events) {\textbf{8. 写事件和状态}\\session event、message、part、runtime state 持久化};
  \node[terminal, below=14mm of events, text width=42mm] (response) {\textbf{9. 返回执行资源}\\blocked / pending request / latest event seq};
  \node[terminal, below=12mm of response, text width=42mm] (sse) {\textbf{10. 前端增量刷新}\\\pathinline{/events/stream}\\+\ \pathinline{/messages}};

  \draw[flow] (user) -- (endpoint);
  \draw[flow] (endpoint) -- (persistuser);
  \draw[signal] (endpoint) -- (versioncheck);
  \draw[flow] (versioncheck.east) -- ++(8mm,0) |- (persistuser.south);
  \draw[flow] (persistuser) -- (loop);
  \draw[flow] (loop) -- (model);
  \draw[flow] (loop) |- (tool);
  \draw[flow] (tool) -- (model);
  \draw[flow] (model) -- (events);
  \draw[flow] (events) -- (response);
  \draw[signal] (events.east) -- ++(10mm,0) |- (sse.east);
  \draw[signal] (sse.west) |- node[pos=0.22, below, font=\scriptsize] {有新事件就继续刷新} (loop.south);

  \node[note, below=12mm of tool, text width=46mm] (toolnote) {只要 session 还存在 pending tool、pending permission、pending user input，\pathinline{run_until_stable()} 就不会直接返回 idle。};
\end{tikzpicture}
\captionof{figure}{Agena 的总主流程。你可以把它理解为一个“先写入用户消息，再循环驱动直到稳定”的状态机。}
\end{diagramlandscape}

\section{运行时启动与热重载流程}

运行时的启动入口来自不同 app，但核心装配动作一致。无论是 CLI、HTTP API server 还是 Studio Server，本质上都会走到 \pathinline{AgenaRuntimeBuilder::build()}。

\subsection{启动阶段在做什么}

\begin{itemize}
  \item 解析 CLI 参数，得到 \pathinline{LoadConfigRequest}。
  \item 连接数据库，必要时创建目录并执行 migration。
  \item 调用 \pathinline{ConfigLoader::load()} 合并配置。
  \item 构建第一代 \pathinline{RuntimeSnapshot}。
  \item 启动 \pathinline{reload::run} 和 \pathinline{janitor::run} 两个后台任务。
\end{itemize}

\subsection{热重载在做什么}

\pathinline{reload::run} 会定期轮询当前 snapshot 的 watch paths。只要发现配置文件或插件目录变化，就重新构建一代新的 snapshot，然后整体替换。

\begin{diagramlandscape}
\begin{tikzpicture}[diagram, node distance=10mm and 14mm, scale=0.92, every node/.style={transform shape}]
  \node[terminal, text width=32mm] (start) {\textbf{程序启动}\\CLI / HTTP API / Studio Server};
  \node[runtime, right=of start, text width=34mm] (args) {\textbf{解析参数}\\config path / mode / overrides / workspace root / database path};
  \node[runtime, right=of args, text width=34mm] (db) {\textbf{连接数据库}\\必要时创建目录\\执行 schema migration};
  \node[runtime, right=of db, text width=34mm] (load) {\textbf{加载配置}\\默认值 + 文件 + mode + 环境变量 + CLI 覆盖};

  \node[runtime, below=18mm of db, text width=40mm] (build) {\textbf{构建 RuntimeSnapshot}\\ProviderRegistry / PluginManager / AuthStore / SessionManager};
  \node[runtime, right=of build, text width=34mm] (serve) {\textbf{对外服务}\\接受 CLI / HTTP / Studio 请求};
  \node[runtime, left=of build, text width=32mm] (janitor) {\textbf{janitor::run}\\定时 prune session cache};

  \node[runtime, right=of serve, text width=34mm] (watch) {\textbf{reload::run}\\周期性读取 watch paths\\比对修改时间戳};
  \node[decision, below=of watch] (changed) {路径有变化？};
  \node[runtime, below=of changed, text width=36mm] (rebuild) {\textbf{重建新 snapshot}\\generation +1\\沿用已有数据库连接};
  \node[runtime, below=of rebuild, text width=36mm] (swap) {\textbf{swap current snapshot}\\新请求读到新配置\\旧快照自然退出};

  \draw[flow] (start) -- (args);
  \draw[flow] (args) -- (db);
  \draw[flow] (db) -- (load);
  \draw[flow] (load.south) -- ++(0,-7mm) -| (build.north);
  \draw[flow] (build) -- (serve);
  \draw[signal] (build) -- (janitor);
  \draw[signal] (build) -- (watch);
  \draw[flow] (watch) -- (changed);
  \draw[flow] (changed) -- node[right, font=\scriptsize] {是} (rebuild);
  \draw[weakflow] (changed.west) -- ++(-12mm,0) node[above, font=\scriptsize] {否} |- (watch.west);
  \draw[flow] (rebuild) -- (swap);
  \draw[flow] (swap.west) -- ++(-12mm,0) |- (serve.south);
\end{tikzpicture}
\captionof{figure}{运行时启动和热重载流程。重点不是“修改局部变量”，而是重建一份新的 RuntimeSnapshot 再整体替换。}
\end{diagramlandscape}

\section{一次用户提问是怎么跑完的}

这里我们只看最典型的一条路径：\highlight{用户在聊天页发送一条新消息}。

\subsection{阶段 1：API 先保证请求是安全可落库的}

HTTP API 在 \pathinline{submit_session_turn()} 里会先做两件事：

\begin{enumerate}
  \item 如果带了 \pathinline{If-Match}，先校验 session version，避免并发覆盖。
  \item 解析 run options。如果请求没显式给 model，就尝试：
  \begin{itemize}
    \item 复用 session 最近一次有效模型；
    \item 如果 runtime 只配置了一个 provider，则退回它的默认模型；
    \item 否则返回 400。
  \end{itemize}
\end{enumerate}

\subsection{阶段 2：先写用户消息，再进入调度循环}

\pathinline{SessionManager::submit_user_turn()} 不会先直接去跑模型，而是先：

\begin{itemize}
  \item 加载 session；
  \item 预留 message ID；
  \item 构建 user message；
  \item 把 user message 持久化；
  \item 然后再进入 \pathinline{run_until_stable()}。
\end{itemize}

这样设计有一个很现实的好处：就算后面的模型调用失败，\highlight{用户消息也已经稳定落库}，前端仍能恢复上下文。

\subsection{阶段 3：SessionManager 不断问“下一步该干什么”}

\pathinline{run_until_stable()} 是 Agena 会话流程里最关键的函数。它每轮都会检查：

\begin{enumerate}
  \item session 是否 blocked；
  \item 是否存在 pending tool；
  \item 当前状态是 idle 还是 awaiting model；
  \item 如果还没稳定，就继续跑下一步。
\end{enumerate}

\begin{notebox}[把它想成一个状态机]
这不是“单次调用模型再返回”的简单双层流程，而是一个多轮循环。模型可以产出工具调用，工具可以触发权限申请或 ask\_user，用户答复后又会继续驱动同一个 session 往前跑，直到回到 idle。
\end{notebox}

\section{模型流、工具流、阻塞流如何衔接}

这一段是整个系统最像 Agent 的地方。

\subsection{模型流}

\pathinline{run_model_turn()} 会：

\begin{itemize}
  \item 先根据 prompt window、上下文预算、tool catalog、provider continuation 能力构造 prompt；
  \item 再让 \pathinline{SessionProcessor::run_turn()} 打开 provider stream；
  \item provider 流式事件通过 \pathinline{MessageProjector} 变成 message / part 的更新；
  \item 流结束后，把 assistant message、usage、provider metadata、event 一起持久化。
\end{itemize}

\subsection{工具流}

如果模型输出了 tool call，那么 session 不会直接 idle，而会出现 pending tool。接下来进入工具处理分支：

\begin{itemize}
  \item 先由 \pathinline{ToolExecutor} 解析 invocation；
  \item 再收集权限检查；
  \item 可能直接执行，也可能插入 permission request part，或者插入 user input request part；
  \item 真正完成后，再生成一条 tool result message；
  \item 然后回到 \pathinline{run_until_stable()}，继续决定下一步是否还要跑模型。
\end{itemize}

\begin{diagramlandscape}
\resizebox{0.94\linewidth}{!}{%
\begin{tikzpicture}[diagram, node distance=10mm and 18mm, scale=0.92, every node/.style={transform shape}]
  \node[runtime] (pendingtool) {\textbf{发现 pending tool}\\tool call 已经写进 assistant message};
  \node[runtime, below=of pendingtool] (prepare) {\textbf{prepare invocation}\\补齐标题、处理 hook、标准化调用};
  \node[decision, below=of prepare] (perm) {需要权限吗？};
  \node[runtime, left=40mm of perm] (askperm) {\textbf{插入 permission request part}\\session 进入 blocked};
  \node[decision, left=40mm of askperm] (replyperm) {用户怎么回复？};
  \node[runtime, below=of perm] (exec) {\textbf{执行工具}\\bash / read / apply\_patch / task / plugin tool};
  \node[decision, below=of exec] (askuser) {工具要求 ask\_user？};
  \node[runtime, right=40mm of askuser] (inputreq) {\textbf{插入 user input request part}\\session 进入 blocked};
  \node[decision, above=of inputreq] (replyinput) {用户是否提交答案？};
  \node[runtime, below=of askuser] (success) {\textbf{写 tool result message}\\记录 loaded deferred tools / attachments / blocks};
  \node[runtime, below=of success] (loopback) {\textbf{回到 run\_until\_stable}\\如果状态仍是 awaiting model，就继续跑模型};
  \node[runtime, left=40mm of success] (failure) {\textbf{写 tool failed}\\reason 进入 tool execution part};

  \draw[flow] (pendingtool) -- (prepare);
  \draw[flow] (prepare) -- (perm);
  \draw[flow] (perm) -- node[above, font=\scriptsize] {是} (askperm);
  \draw[flow] (askperm) -- (replyperm);
  \draw[flow] (replyperm) |- node[pos=0.25, left, font=\scriptsize] {允许} (exec);
  \draw[flow] (replyperm) |- node[pos=0.25, above, font=\scriptsize] {拒绝} (failure);
  \draw[flow] (perm) -- node[right, font=\scriptsize] {否} (exec);
  \draw[flow] (exec) -- (askuser);
  \draw[flow] (askuser) -- node[above, font=\scriptsize] {是} (inputreq);
  \draw[flow] (inputreq) -- (replyinput);
  \draw[flow] (replyinput) |- node[pos=0.25, above, font=\scriptsize] {提交} (success);
  \draw[flow] (replyinput) |- node[pos=0.2, right, font=\scriptsize] {取消} (failure);
  \draw[flow] (askuser) -- node[right, font=\scriptsize] {否} (success);
  \draw[flow] (success) -- (loopback);
  \draw[flow] (failure) -- (loopback);
\end{tikzpicture}
}
\captionof{figure}{工具调用分支流程。Agena 把“权限阻塞”和“用户补充输入阻塞”都显式建模成 part，因此 session 可以在中途停下、恢复、再继续。}
\end{diagramlandscape}

\section{为什么前端能实时看到变化}

前端不是每次都把整段聊天历史整包重拉，而是结合了三种机制：

\begin{itemize}
  \item \pathinline{GET /sessions/\{id\}/state}
  \begin{itemize}
    \item 用来恢复 blocked、pending request、latest event seq。
  \end{itemize}
  \item \pathinline{GET /sessions/\{id\}/events/stream}
  \begin{itemize}
    \item 用 SSE 按 seq 增量推送 session event。
  \end{itemize}
  \item \pathinline{GET /sessions/\{id\}/messages?parts=summary}
  \begin{itemize}
    \item 只拉 message summary，真正需要 detail 再单独取。
  \end{itemize}
\end{itemize}

从前端 \pathinline{ChatPage.vue} 可以看出，它会优先尝试 SSE；如果运行环境不支持，就回退到定时轮询。

\begin{diagramlandscape}
\begin{tikzpicture}[diagram, node distance=14mm and 18mm]
  \node[terminal] (chat) {\textbf{ChatPage.vue}\\选中 session 后开始恢复会话};
  \node[terminal, right=of chat] (state) {\textbf{GET session state}\\拿 blocked / run\_state / latest\_event\_seq};
  \node[terminal, right=of state] (messages) {\textbf{GET messages summary}\\按 created\_at 正序展示};
  \node[terminal, right=of messages] (sseopen) {\textbf{打开 SSE}\\\pathinline{/events/stream?after\_seq=n}};

  \node[decision, below=of sseopen] (support) {环境支持\\ReadableStream 吗？};
  \node[terminal, left=38mm of support] (poll) {\textbf{轮询 fallback}\\定时刷新 state + messages};
  \node[terminal, below=of support] (event) {\textbf{收到 session event}\\更新 latest seq\\本地先做轻量 patch};
  \node[terminal, below=of event] (refresh) {\textbf{调度刷新}\\短延迟后补拉 messages / state};
  \node[terminal, left=38mm of refresh] (render) {\textbf{重新渲染 UI}\\显示文本块、tool part、permission / ask\_user 卡片};

  \draw[flow] (chat) -- (state);
  \draw[flow] (state) -- (messages);
  \draw[flow] (messages) -- (sseopen);
  \draw[flow] (sseopen) -- (support);
  \draw[flow] (support) -- node[above, font=\scriptsize] {否} (poll);
  \draw[flow] (support) -- node[right, font=\scriptsize] {是} (event);
  \draw[signal] (poll) |- (refresh);
  \draw[signal] (event) -- (refresh);
  \draw[flow] (refresh) -- (render);
  \draw[weakflow] (render.west) -- ++(-12mm,0) |- node[pos=0.25, left, font=\scriptsize] {继续等待新事件} (support.west);
\end{tikzpicture}
\captionof{figure}{前端增量刷新流程。Agena 不是只靠一条 SSE 就把所有内容推完，而是用 SSE 指挥“什么时候需要刷新”，再结合摘要接口做局部 hydration。}
\end{diagramlandscape}

\section{Provider 鉴权流程怎么走}

从代码看，Agena 的鉴权不是只支持一种方式，而是至少覆盖了三类：

\begin{itemize}
  \item Browser OAuth：例如 OpenAI、GitLab。
  \item Device Code：例如 OpenAI 设备码、GitHub Copilot 设备码。
  \item API Key：例如 OpenAI / Anthropic / OpenAI-compatible。
\end{itemize}

\pathinline{AuthManager} 负责统一这些登录流程，并最终把凭据写进 \pathinline{AuthStore}。默认实现是文件存储，也就是本地的 auth data 文件。

\begin{diagramlandscape}
\begin{tikzpicture}[diagram, node distance=12mm and 16mm]
  \node[terminal] (client) {\textbf{前端 / CLI / 操作员}\\发起登录或写入 API key};
  \node[terminal, right=of client] (authapi) {\textbf{Auth API}\\\pathinline{/auth/providers/*}\\start / finish / poll / refresh / api-key};
  \node[runtime, right=of authapi] (manager) {\textbf{AuthManager}\\统一封装 OpenAI / GitLab / Copilot 登录逻辑};

  \node[runtime, above right=12mm and 18mm of manager] (browseroauth) {\textbf{Browser OAuth 分支}\\start 返回 authorize URL + state + pkce\\finish 交换 token};
  \node[runtime, below=14mm of manager] (apikey) {\textbf{API Key 分支}\\直接校验非空并写入 store};
  \node[runtime, below right=12mm and 18mm of manager] (devicecode) {\textbf{Device Code 分支}\\start 返回 verification URL + user code\\poll 等待完成};

  \node[external, right=22mm of manager] (providerauth) {\textbf{Provider 鉴权服务}\\OpenAI / GitLab / GitHub 等授权端点};
  \node[storage, below=14mm of providerauth] (store) {\textbf{AuthStore}\\默认是文件存储\\保存 API key / access token / refresh token / expires};

  \draw[flow] (client) -- (authapi);
  \draw[flow] (authapi) -- (manager);
  \draw[flow] (manager) -- (browseroauth);
  \draw[flow] (manager) -- (devicecode);
  \draw[flow] (manager) -- (apikey);
  \draw[flow] (browseroauth) -- (providerauth);
  \draw[flow] (devicecode) -- (providerauth);
  \draw[flow] (browseroauth) |- (store);
  \draw[flow] (devicecode) |- (store);
  \draw[flow] (apikey.east) -| (store.west);

  \node[note, below=10mm of store] (authnote) {auth 读取接口只返回安全摘要，不会把明文 secret 回传给前端。};
\end{tikzpicture}
\captionof{figure}{Provider 鉴权流程。无论是 OAuth、Device Code 还是 API Key，最终都会被归一到同一个 AuthStore。}
\end{diagramlandscape}

\section{常见阻塞点与恢复点}

\begin{longtable}{p{0.22\textwidth}p{0.31\textwidth}p{0.33\textwidth}}
\toprule
\textbf{阻塞点} & \textbf{为什么会停住} & \textbf{如何恢复} \\
\midrule
\endhead
Permission request & 工具命中了 ask 型权限策略，或持久化规则要求确认。 & 前端调用 \pathinline{/permission-replies}，回复 Allow/Deny。 \\
User input request & 工具执行过程中需要向用户补问题。 & 前端调用 \pathinline{/user-input-replies}，提交答案或取消。 \\
Provider stream error & 上游模型接口返回失败、超时、流中断。 & 查看 session event 和错误信息，必要时修复配置后继续提交。 \\
Version conflict & 并发修改 session，\pathinline{If-Match} 对不上。 & 客户端重新拉最新 session version，再重试写操作。 \\
Reload 生效前后的配置差异 & watch path 触发了新的 runtime snapshot。 & 前端重新拉 \pathinline{/runtime} 和 provider 列表，按新 snapshot 行为工作。 \\
\bottomrule
\end{longtable}

\section{本篇总结}

如果把 Agena 的流程压缩成一句最朴素的话，可以这样说：

\begin{summarybox}[最终心智模型]
\term{Agena} 会先把用户动作\highlight{稳定写入}，然后用一个 \highlight{run\_until\_stable} 循环不断驱动 session 往前走。这个循环可以跑模型、执行工具、插入权限请求、插入用户补充输入请求，也可以在任意一步停住并在下一次请求中恢复。前端则通过 \pathinline{session state + SSE + messages summary} 跟上这个过程。
\end{summarybox}

到这里，你应该已经能把“架构”和“流程”连在一起看了：

\begin{itemize}
  \item 运行时为什么要快照化；
  \item 会话为什么要拆成 manager / processor / store / executor；
  \item 前端为什么既要 SSE 又要摘要接口；
  \item 权限阻塞和用户补充输入为什么都建模成 part；
  \item 为什么 Studio 和 HTTP API 可以共享同一套资源语义。
\end{itemize}

\end{document}
